% Canonical LaTeX chapter; migrated from docs/pack-format.md.
\hypertarget{pack-format}{%
\subsection{Pack format}\label{pack-format}}

Packs are UTF-8 TOML parsed by Python 3.11+ \texttt{tomllib}. See
\href{../packs/community/}{\texttt{packs/community/}} for complete
examples in cyber, Python, and 3D art. The examples are drafts for human
review, not generated lesson content.

\begin{verbatim}
[pack]
id = "tiny-tool"
title = "Build a tiny file tool"
kind = "project"
version = 1
source = "human"
reviewed = false

[vars]
work_dir = "~/learning"

[[level]]
id = "first"
title = "A working example"
lane = "main"
gate = true

[[node]]
id = "example"
level = "first"
title = "Save a working example"
size = "S"
minutes = 20
needs = []
done_when = "A working example saved with an explanation in my own words"
preview = "Build a tool for organizing a small set of files."
requirements = "Python 3 and a text editor."
brief = "Build a small tool that solves one annoying file-management problem."
steps = ["Choose a tiny input.", "Run the tool and compare its output with your prediction."]
hints = ["Start with one file before handling a whole directory."]
resource = "https://docs.python.org/3/tutorial/"
verify = {type = "file", path = "{work_dir}/example.py", min_bytes = 20}
\end{verbatim}

Required pack fields: kebab-case \texttt{id}, \texttt{title},
\texttt{kind} (\texttt{project} or \texttt{proficiency}), and positive
integer \texttt{version}. \texttt{source} is \texttt{human},
\texttt{ai-draft}, or \texttt{imported}. Every node has an existing
\texttt{level}, a unique stable \texttt{id}, a \texttt{title}, S/M/L
\texttt{size}, a \texttt{done\_when} of at least 15 characters, and an
HTTP(S) \texttt{resource}.

Optional \texttt{brief} text introduces a concrete challenge. Optional
\texttt{steps} and \texttt{hints} are lists of nonempty strings, with
the same variable and series substitutions as other content. The web
started challenge workspace displays steps in order and keeps hints behind
individual disclosures. These are learning guidance, not completion
checks; \texttt{done\_when} and the verifier still define the finish
line. Older packs without them continue to work. Prefer a small
observable result over a reading assignment, and set honest
\texttt{minutes} estimates so short sessions have useful tasks.

Optional \texttt{preview} states a short objective before Start.
\texttt{requirements} lists prerequisite tools, access, hardware and costs;
exclude task files, commands, answers and detailed completion criteria. An older
pack without these fields uses its title as the preview. Optional \texttt{setup}
contains workspace preparation instructions displayed after Start. \texttt{resource\_label} and \texttt{resource\_note} describe a
reference's role. The resource link, label and note appear before Start and remain
available in the workspace; opening a resource does not start timing.
\texttt{materials} contains label/URL objects for actual inputs;
HTTP(S) and bounded \path{/training/} ZIP/JSON/CSV/TXT paths are accepted. Fenced
code in steps is copyable text, never an automatically executed command. Older
packs omit these fields. See the challenge-guidance implementation chapter for
the generator, downloads and evidence workflow.

S/M/L conservatively fit 30/120/240-minute budgets. Optional integer
\texttt{minutes} may refine the estimate within that size bound.
\texttt{needs} is a list of node IDs forming a directed acyclic graph.
\texttt{series = \{from = 13, to = 15\}} expands an ID
\texttt{exercise} to \texttt{exercise-13}, \texttt{exercise-14},
\texttt{exercise-15}, chained in order. \texttt{\{n\}} and
\texttt{\{next\}} are substituted. The first expanded node inherits the
original needs.

Main levels gate on the previous main level. Milestones
(\texttt{milestone = true}) replace the all-node gate when present.
Proficiency packs require 8--15 main levels and at least one M/L
milestone per main level. Side levels are always open but individual
prerequisites still apply. \texttt{gate = false} bypasses a level
gate. \texttt{ascend unlock \textless{}level\textgreater{}} records an
explicit early unlock.

\texttt{{[}vars{]}} values substitute \texttt{\{name\}} placeholders;
local verifier paths expand \texttt{\textasciitilde{}}. For commands,
the template is split into arguments \textbf{before} substitution. Shell
operators are not interpreted. A command containing pipelines or
redirects must be rewritten as an executable script and trusted
deliberately.

Install with
\texttt{ascend pack install \textless{}file-or-bundled-id\textgreater{}}.
Private copies live under
\texttt{ASCEND\_HOME/packs/\textless{}user-id\textgreater{}/}. Editing a
copy resets trust and review state.
\texttt{ascend pack diff \textless{}updated-file\textgreater{}}
reports added/removed IDs; progress is retained for stable IDs.
\texttt{ascend pack review \textless{}id\textgreater{}} records
review without editing the file.

Available checks: \texttt{manual}, \texttt{file}, \texttt{command},
\texttt{http}, \texttt{git}, and
\texttt{plugin:\textless{}name\textgreater{}}.
\texttt{auto\_complete = true} permits passing hard local file/command
checks to complete without a separate confirmation.
\texttt{ascend sync} proposes other unlocked results;
\texttt{ascend accept} reruns the checks before accepting. Known
evidence timestamps produce deduplicated entries on their actual logical
day, unless \texttt{ascend config auto-counts-for-streak false} is
set.

\texttt{ascend split \textless{}node\textgreater{} "first small result" "second small result"}
stores a local 10-minute-step breakdown in SQLite; the TOML input stays
untouched. Choose parts that together satisfy the
parent\textquotesingle s original completion condition.


\subsubsection{Authored checkpoints}
A node may include up to twelve \texttt{checkpoints}, each an inline table with
\texttt{id}, \texttt{title}, \texttt{minutes}, \texttt{done\_when} and
\texttt{next\_step}. IDs must be unique within the node; minutes must be positive
and shorter than the full task estimate. Author meaningful stopping points, not
cosmetic fragments or shorter labels for a complete assessment. Saving one records
a partial learning note and does not replace the parent's verification or unlock
later levels. Checkpoint IDs remain stable across updates. Existing installed packs
are not overwritten when bundled content gains checkpoints.

\begin{verbatim}
checkpoints = [{ id = "predict", title = "Predict one result",
  minutes = 10, done_when = "Compare a prediction with one observation.",
  next_step = "Complete the remaining cases in the parent task." }]
\end{verbatim}
In TOML, write an inline table on one line (the example is wrapped for print).

\subsubsection{Playground content is separate from track nodes}
Repeating optional TOML nodes does not create recurring practice. Playground's
initial bundled catalog is authored in \path{ascend/playground/catalog.py}, with
fixtures and evaluators in \path{ascend/playground/fixtures.py} and subject-specific
content in \path{ascend/playground/subjects.py}. Example steps and reflection prompts
live in \path{ascend/playground/guidance.py}. This release does
not import activity definitions from third-party TOML packs. Additional subjects
can extend that catalog without mission enrollment or changing stable task IDs.

Each definition needs a version, independent activity ID, kind, example/challenge
role, skill abilities, preparation references and time, honest working/setup
estimates, device/environment needs, evidence mechanism, variant strategy, stable
reward key and optional source-node identity. Reused assessment identities must
keep their reward key; changing a title or version must not mint another first-solve
award. Source-node references must also retain the original assessment reward.
The initial Playground challenges use distinct fixtures from track assessments.

Use neutral mystery IDs/titles and broad prerequisite abilities. Keep expected
answers, diagnosis, hints and debriefs outside public discovery metadata. Author
observable evidence, plausible contrasting or benign cases, a discriminating test
or prediction, and a meaningful changed-condition follow-up. Free text alone is
not automatically verified reasoning. Examples deliberately teach their own
separate fixture. A challenge may have no hints or limited authored hints; it does
not need a guided conversion. Finite variants must be labeled review after they
are seen, and a solution reveal must remain associated with the same fixture.
